% GENERATED FILE. Edit canonical lessons, then run npm run generate:books.
% canonical-source-hash: fnv1a64:2169a95c984f8b21
% canonical-lessons: ML-C219-nanayam, ML-C219-nottu, ML-C219-cillara, ML-C219-bakki, ML-C219-kilivu

\chapter{Coin, Banknote, Small change, Balance, Discount}
\label{ch:ml-coin-banknote-small-change-balance-discount}
\hlchaptermodality{219}

\begin{chapteropening}
\textbf{By the end of this chapter:} \emph{I can say a coin, a banknote, small change, the balance and a discount.}

Complete the last lesson of chapter 219: I can say a coin, a banknote, small change, the balance and a discount.
\end{chapteropening}

\section[\texorpdfstring{\ml{നാണയം} (nāṇayaṁ)}{നാണയം (nāṇayaṁ)}]{\ml{നാണയം} (nāṇayaṁ) — a coin}

\label{lesson:ML-C219-nanayam}



\begin{quote}
Before the new one: say the Malayalam for elephant-foot yam, then the Malayalam for a pumpkin.
\end{quote}

\subsection*{You'll want to know: \ml{നാണയം}}
\textbf{\ml{നാണയം}} — \emph{nāṇayaṁ} — \textquotedblleft{}a coin\textquotedblright{}.

\begin{grammarlens}[title={What you've built}]
One more word for this chapter.
\end{grammarlens}

\subsection*{Your turn}
\begin{itemize}
\raggedright
  \item \emph{Say it:} \emph{nāṇayaṁ}
  \item \emph{Say it:} \emph{nāṇayaṁ}, once more
  \item \emph{Say it:} say \emph{nāṇayaṁ}
  \item \emph{Recall:} say the Malayalam for elephant-foot yam, then the Malayalam for a pumpkin, then say \emph{nāṇayaṁ} again
\end{itemize}

\begin{tcolorbox}[breakable,colback=teal!4,colframe=teal!35!black,title={Before you move on}]
What does \ml{നാണയം} mean? (\textquotedblleft{}A coin\textquotedblright{}.) Say it once more.
\end{tcolorbox}
\section[\texorpdfstring{\ml{നോട്ട്} (nōṭṭŭ)}{നോട്ട് (nōṭṭŭ)}]{\ml{നോട്ട്} (nōṭṭŭ) — a banknote}

\label{lesson:ML-C219-nottu}



\begin{quote}
Before the new one: say the Malayalam for a pumpkin, then the Malayalam for a coin.
\end{quote}

\subsection*{You'll want to know: \ml{നോട്ട്}}
\textbf{\ml{നോട്ട്}} — \emph{nōṭṭŭ} — \textquotedblleft{}a banknote\textquotedblright{}.

\begin{grammarlens}[title={What you've built}]
One more word for this chapter.
\end{grammarlens}

\subsection*{Your turn}
\begin{itemize}
\raggedright
  \item \emph{Say it:} \emph{nōṭṭŭ}
  \item \emph{Say it:} \emph{nōṭṭŭ}, once more
  \item \emph{Say it:} say \emph{nōṭṭŭ}
  \item \emph{Recall:} say the Malayalam for a pumpkin, then the Malayalam for a coin, then say \emph{nōṭṭŭ} again
\end{itemize}

\begin{tcolorbox}[breakable,colback=teal!4,colframe=teal!35!black,title={Before you move on}]
What does \ml{നോട്ട്} mean? (\textquotedblleft{}A banknote\textquotedblright{}.) Say it once more.
\end{tcolorbox}
\section[\texorpdfstring{\ml{ചില്ലറ} (cillaṟa)}{ചില്ലറ (cillaṟa)}]{\ml{ചില്ലറ} (cillaṟa) — small change (coins)}

\label{lesson:ML-C219-cillara}



\begin{quote}
Before the new one: say the Malayalam for a coin, then the Malayalam for a banknote.
\end{quote}

\subsection*{You'll want to know: \ml{ചില്ലറ}}
\textbf{\ml{ചില്ലറ}} — \emph{cillaṟa} — \textquotedblleft{}small change (coins)\textquotedblright{}.

\begin{grammarlens}[title={What you've built}]
One more word for this chapter.
\end{grammarlens}

\subsection*{Your turn}
\begin{itemize}
\raggedright
  \item \emph{Say it:} \emph{cillaṟa}
  \item \emph{Say it:} \emph{cillaṟa}, once more
  \item \emph{Say it:} say \emph{cillaṟa}
  \item \emph{Recall:} say the Malayalam for a coin, then the Malayalam for a banknote, then say \emph{cillaṟa} again
\end{itemize}

\begin{tcolorbox}[breakable,colback=teal!4,colframe=teal!35!black,title={Before you move on}]
What does \ml{ചില്ലറ} mean? (\textquotedblleft{}Small change (coins)\textquotedblright{}.) Say it once more.
\end{tcolorbox}
\section[\texorpdfstring{\ml{ബാക്കി} (bākki)}{ബാക്കി (bākki)}]{\ml{ബാക്കി} (bākki) — the balance, the change given back}

\label{lesson:ML-C219-bakki}



\begin{quote}
Before the new one: say the Malayalam for a banknote, then the Malayalam for small change.
\end{quote}

\subsection*{You'll want to know: \ml{ബാക്കി}}
\textbf{\ml{ബാക്കി}} — \emph{bākki} — \textquotedblleft{}the balance, the change given back\textquotedblright{}.

\begin{grammarlens}[title={What you've built}]
One more word for this chapter.
\end{grammarlens}

\subsection*{Your turn}
\begin{itemize}
\raggedright
  \item \emph{Say it:} \emph{bākki}
  \item \emph{Say it:} \emph{bākki}, once more
  \item \emph{Say it:} say \emph{bākki}
  \item \emph{Recall:} say the Malayalam for a banknote, then the Malayalam for small change, then say \emph{bākki} again
\end{itemize}

\begin{tcolorbox}[breakable,colback=teal!4,colframe=teal!35!black,title={Before you move on}]
What does \ml{ബാക്കി} mean? (\textquotedblleft{}The balance, the change given back\textquotedblright{}.) Say it once more.
\end{tcolorbox}
\section[\texorpdfstring{\ml{കിഴിവ്} (kiḻivŭ)}{കിഴിവ് (kiḻivŭ)}]{\ml{കിഴിവ്} (kiḻivŭ) — a discount}

\label{lesson:ML-C219-kilivu}



\begin{quote}
Before the new one: say the Malayalam for small change, then the Malayalam for the balance.
\end{quote}

\subsection*{You'll want to know: \ml{കിഴിവ്}}
\textbf{\ml{കിഴിവ്}} — \emph{kiḻivŭ} — \textquotedblleft{}a discount\textquotedblright{}.

\begin{grammarlens}[title={What you've built}]
One more word for this chapter.
\end{grammarlens}

\subsection*{Your turn}
\begin{itemize}
\raggedright
  \item \emph{Say it:} \emph{kiḻivŭ}
  \item \emph{Say it:} \emph{kiḻivŭ}, once more
  \item \emph{Say it:} say \emph{kiḻivŭ}
  \item \emph{Recall:} say the Malayalam for small change, then the Malayalam for the balance, then say \emph{kiḻivŭ} again
\end{itemize}

\begin{tcolorbox}[breakable,colback=teal!4,colframe=teal!35!black,title={Before you move on}]
What does \ml{കിഴിവ്} mean? (\textquotedblleft{}A discount\textquotedblright{}.) Say it once more.
\end{tcolorbox}
